\section{Event Scheduling and Numerical Fidelity}
\subsection{Notification replaces failed admission attempts}
Periodic polling can make simulator execution cost depend on how long an unavailable resource remains unavailable. If $K$ blocked producers each retry every $\Delta$ time units for duration $T$, their failed checks can scale approximately as $KT/\Delta$. These checks do not transfer useful data. An event-driven adapter instead registers waiting producers and schedules a retry when a relevant resource changes state. This reduces simulator work, but it must preserve the resource constraints rather than granting capacity on notification.

The current IO adapter keeps waiting producers associated with their original retry-clock phase. When queue space becomes available, it computes the next eligible retry time for a waiting phase and arms a port-level event. Within that phase, registered producers are resumed in registration order while space remains. Device egress blockage uses a release notification from the endpoint device. A notification is an opportunity to retry admission, not proof that a complete packet can be transmitted without checking egress and credit state.

This scheduling change does not alter AFH forwarding or the LLR/CBFC protocol algorithms. It can nevertheless change the order in which simultaneous callbacks enter an adapter queue. Since packet formation depends on which compatible records are present at that instant, different event insertion orders can change packet sizes and individual task times. We therefore distinguish protocol invariance, resource-contract invariance, and trace equivalence. The first two do not establish the third.

For an optimization claimed to preserve a particular experiment, we compare its input hashes, total useful bytes, completion status, and timing records. If traces differ, we report the difference instead of treating delivery alone as proof of timing equivalence. The PCIe case in this paper passes a stronger paired check: four families of output files match exactly between wakeup modes at each precision. That observation is deliberately scoped to the paired workload and cannot be generalized to all queue-contention patterns.

\subsection{Clock resolution is not hardware issue frequency}
A discrete-event simulator advances to scheduled events rather than executing one iteration for every time quantum. Moving from 1\,ns to 1\,ps resolution therefore does not imply one thousand times as many event callbacks. The event count depends on modeled operations and scheduling decisions. A 200\,ns memory access remains a delayed event; a 1\,ns admission interval remains a 1\,ns interval. Finer timestamps allow the model to represent service and propagation durations between those boundaries.

There are at least three timing contracts: event timestamp representation, resource-service granularity, and producer issue frequency. The current configuration separates them. Global time is represented in picoseconds. MPU bandwidth service proceeds in 1\,ns windows with a byte budget. IO producers retain their configured request intervals, and packet formation has its own opportunity schedule. A future cycle-accurate IO model would additionally constrain how many operations a particular hardware block can accept each cycle; timestamp precision alone supplies no such constraint.

When a byte server has a fixed cycle, it should conserve residual capacity within the cycle. Serving a 256-byte transaction on a 190-byte-per-cycle server requires at least two service windows for that transaction, but the second window need not be dedicated exclusively to its remaining 66 bytes. The next ready FIFO transaction may consume the remaining 124 bytes. Rounding every transaction to two exclusive cycles would create a different and lower-bandwidth server. The same principle explains why coarse per-packet serialization can distort the PCIe reference path.

\subsection{Invariants and reproducibility boundaries}
We use three classes of checks. Structural checks reject invalid dependency edges, illegal endpoint identities, and inconsistent compiled inputs. Conservation checks compare generated, transmitted, received, and resource-serviced bytes while distinguishing logical data from protocol overhead. Temporal checks verify that a consumer starts only after its declared completion condition and that resource occupancy and service order match the selected model.

These checks establish internal consistency; none substitutes for external validation. A precisely reproducible model can be consistently wrong about a switch pipeline or an IO Die. Conversely, measurements can reveal the correct throughput while leaving insufficient evidence to explain queue occupancy. We retain both functional invariants and a separate calibration requirement so that a successful regression is not mislabeled as silicon accuracy.

An experiment records resolved inputs, component hashes, the simulator binary identity, status, and raw statistics. Multiple conditions execute as independent processes; this is different from parallelizing one simulated timeline. Capacity-rejected conditions are reported as such, without assigning a completion time. Timeouts and unfinished graphs remain failures rather than disappearing from an average. The harness facilitates these practices, but its presence is not itself evidence for the scientific claims.
